%Book
%Copyright (C) 2019  Patrick Diehl
%
%This program is free software: you can redistribute it and/or modify
%it under the terms of the GNU General Public License as published by
%the Free Software Foundation, either version 3 of the License, or
%(at your option) any later version.

%This program is distributed in the hope that it will be useful,
%but WITHOUT ANY WARRANTY; without even the implied warranty of
%MERCHANTABILITY or FITNESS FOR A PARTICULAR PURPOSE.  See the
%GNU General Public License for more details.

%You should have received a copy of the GNU General Public License
%along with this program.  If not, see <http://www.gnu.org/licenses/>.

% The template of the book is based on the The Legrand Orange Book
% LaTeX Template
% Version 2.4 (26/09/2018)
% Original author:
% Mathias Legrand (legrand.mathias@gmail.com) with modifications by:
% Vel (vel@latextemplates.com)
% License:
% CC BY-NC-SA 3.0 (http://creativecommons.org/licenses/by-nc-sa/3.0/)

%----------------------------------------------------------------------------------------
%	PACKAGES AND OTHER DOCUMENT CONFIGURATIONS
%----------------------------------------------------------------------------------------

\documentclass[11pt,fleqn]{book} % Default font size and left-justified equations

%%%%%%%%%%%%%%%%%%%%%%%%%%%%%%%%%%%%%%%%%
% The Legrand Orange Book
% Structural Definitions File
% Version 2.1 (26/09/2018)
%
% Original author:
% Mathias Legrand (legrand.mathias@gmail.com) with modifications by:
% Vel (vel@latextemplates.com)
% 
% This file was downloaded from:
% http://www.LaTeXTemplates.com
%
% License:
% CC BY-NC-SA 3.0 (http://creativecommons.org/licenses/by-nc-sa/3.0/)
%
%%%%%%%%%%%%%%%%%%%%%%%%%%%%%%%%%%%%%%%%%

%----------------------------------------------------------------------------------------
%	VARIOUS REQUIRED PACKAGES AND CONFIGURATIONS
%----------------------------------------------------------------------------------------
\usepackage{fontspec}
\setmainfont{Raleway}

\usepackage[
    type={CC},
    modifier={by-nc-nd},
    version={4.0},
]{doclicense} 

\usepackage{graphicx} % Required for including pictures
\graphicspath{{Pictures/}} % Specifies the directory where pictures are stored
\usepackage{subcaption}

\usepackage{lipsum} % Inserts dummy text


\usepackage{tikz} % Required for drawing custom shapes
\usetikzlibrary{mindmap}
\usetikzlibrary{shapes.geometric, arrows,matrix}
\usetikzlibrary{calc}
\usetikzlibrary{fit}
\tikzstyle{startstop} = [rectangle, rounded corners, minimum width=3cm, minimum height=0.5cm,text centered, draw=black, fill=cadetgrey!30]
\tikzstyle{io} = [trapezium, trapezium left angle=70, trapezium right angle=110, minimum width=3cm, minimum height=0.5cm, text centered, draw=black] %, fill=blue!30]
\tikzstyle{process} = [rectangle, minimum width=3cm, minimum height=0.5cm, text centered, draw=black] %, fill=orange!30]
\tikzstyle{decision} = [diamond, minimum width=2cm, minimum height=0.25cm, text centered, draw=black] %, fill=green!30]
\tikzstyle{arrow} = [thick,->,>=stealth]
\usepackage{pgfplots}

\newcommand\montecarlo[2]{
\def\x{0}
\def\y{0}
\def\k{0}
\def\n{0}
\def\r{0}
\def\pi{3.1415}
\def\e{0}
    \foreach \i in {1,2,...,#1}{%
        \typeout{Point \i}%
        \pgfmathsetmacro\x{rnd}%
        \typeout{X \x}%
        \pgfmathsetmacro\y{rnd}%
        \typeout{Y \y}%
        \pgfmathsetmacro\k{(pow(\x,2)+pow(\y,2)) <pow(1,2)}%
        \typeout{im Kreis?: \k}%
        \pgfmathparse{ifthenelse(\k==1,1,0)}%
        \pgfmathparse{\n+\pgfmathresult}
        \xdef\n{\pgfmathresult}
        \typeout{N \n} 
        }
\pgfmathparse{4*\n/#1}
\xdef\r{\pgfmathresult}
\typeout{R \r}
\pgfmathparse{(\r-\pi)/\pi*100}
\xdef\e{abs(\pgfmathresult)}
\typeout{E \e}
\addplot [only marks,mark=*] coordinates { (#2,\e) };		        
}

\usepackage{endnotes}

\usepackage[english]{babel} % English language/hyphenation

\usepackage{enumitem} % Customize lists
\setlist{nolistsep} % Reduce spacing between bullet points and numbered lists

\usepackage{booktabs} % Required for nicer horizontal rules in tables

\usepackage{xcolor} % Required for specifying colors by name
\definecolor{azure}{rgb}{0.0, 0.5, 1.0}
\definecolor{darkgreen}{rgb}{0.0, 0.5, 0.0}
\definecolor{amaranth}{rgb}{0.9, 0.17, 0.31}
\definecolor{cadetgrey}{rgb}{0.57, 0.64, 0.69}


\usepackage{listings}

\lstset{
    basicstyle=\ttfamily,
    keywordstyle=\color{azure}\ttfamily,
    stringstyle=\color{amaranth}\ttfamily,
    commentstyle=\color{darkgreen}\ttfamily,
    morecomment=[l][\color{magenta}]{\#},
    breaklines=true,
    numbers=left,
    numberstyle=\tiny,
    frame = single,
}

\usepackage{xfrac}

\newcommand{\orcid}[1]{\href{https://orcid.org/#1}{\includegraphics[height=10pt]{orcid}}}
\newcommand{\link}[1]{\endnote{\url{#1}}}
\newcommand{\cpp}[1]{\lstinline[language=c++]{#1}}

%----------------------------------------------------------------------------------------
%	MARGINS
%----------------------------------------------------------------------------------------

\usepackage{geometry} % Required for adjusting page dimensions and margins

\geometry{
	paper=a4paper, % Paper size, change to letterpaper for US letter size
	top=3cm, % Top margin
	bottom=3cm, % Bottom margin
	left=3cm, % Left margin
	right=3cm, % Right margin
	headheight=14pt, % Header height
	footskip=1.4cm, % Space from the bottom margin to the baseline of the footer
	headsep=10pt, % Space from the top margin to the baseline of the header
	%showframe, % Uncomment to show how the type block is set on the page
}

%----------------------------------------------------------------------------------------
%	FONTS
%----------------------------------------------------------------------------------------

\usepackage{avant} % Use the Avantgarde font for headings
%\usepackage{times} % Use the Times font for headings
\usepackage{mathptmx} % Use the Adobe Times Roman as the default text font together with math symbols from the Sym­bol, Chancery and Com­puter Modern fonts

\usepackage{microtype} % Slightly tweak font spacing for aesthetics
\usepackage[utf8]{inputenc} % Required for including letters with accents
\usepackage[T1]{fontenc} % Use 8-bit encoding that has 256 glyphs

%----------------------------------------------------------------------------------------
%	BIBLIOGRAPHY AND INDEX
%----------------------------------------------------------------------------------------

\usepackage[style=numeric,citestyle=numeric,sorting=nyt,sortcites=true,autopunct=true,babel=hyphen,hyperref=true,abbreviate=false,backref=true,backend=biber]{biblatex}
\addbibresource{bib.bib} % BibTeX bibliography file
\defbibheading{bibempty}{}

\usepackage{calc} % For simpler calculation - used for spacing the index letter headings correctly
\usepackage{imakeidx} % Required to make an index
\makeindex[options=-s index.ist] % Tells LaTeX to create the files required for indexing

%----------------------------------------------------------------------------------------
%	MAIN TABLE OF CONTENTS
%----------------------------------------------------------------------------------------

\usepackage{titletoc} % Required for manipulating the table of contents

\contentsmargin{0cm} % Removes the default margin

% Part text styling (this is mostly taken care of in the PART HEADINGS section of this file)
\titlecontents{part}
	[0cm] % Left indentation
	{\addvspace{20pt}\bfseries} % Spacing and font options for parts
	{}
	{}
	{}

% Chapter text styling
\titlecontents{chapter}
	[1.25cm] % Left indentation
	{\addvspace{12pt}\large\sffamily\bfseries} % Spacing and font options for chapters
	{\color{azure!60}\contentslabel[\Large\thecontentslabel]{1.25cm}\color{azure}} % Formatting of numbered sections of this type
	{\color{azure}} % Formatting of numberless sections of this type
	{\color{azure!60}\normalsize\;\titlerule*[.5pc]{.}\;\thecontentspage} % Formatting of the filler to the right of the heading and the page number

% Section text styling
\titlecontents{section}
	[1.25cm] % Left indentation
	{\addvspace{3pt}\sffamily\bfseries} % Spacing and font options for sections
	{\contentslabel[\thecontentslabel]{1.25cm}} % Formatting of numbered sections of this type
	{} % Formatting of numberless sections of this type
	{\hfill\color{black}\thecontentspage} % Formatting of the filler to the right of the heading and the page number

% Subsection text styling
\titlecontents{subsection}
	[1.25cm] % Left indentation
	{\addvspace{1pt}\sffamily\small} % Spacing and font options for subsections
	{\contentslabel[\thecontentslabel]{1.25cm}} % Formatting of numbered sections of this type
	{} % Formatting of numberless sections of this type
	{\ \titlerule*[.5pc]{.}\;\thecontentspage} % Formatting of the filler to the right of the heading and the page number

% Figure text styling
\titlecontents{figure}
	[1.25cm] % Left indentation
	{\addvspace{1pt}\sffamily\small} % Spacing and font options for figures
	{\thecontentslabel\hspace*{1em}} % Formatting of numbered sections of this type
	{} % Formatting of numberless sections of this type
	{\ \titlerule*[.5pc]{.}\;\thecontentspage} % Formatting of the filler to the right of the heading and the page number

% Table text styling
\titlecontents{table}
	[1.25cm] % Left indentation
	{\addvspace{1pt}\sffamily\small} % Spacing and font options for tables
	{\thecontentslabel\hspace*{1em}} % Formatting of numbered sections of this type
	{} % Formatting of numberless sections of this type
	{\ \titlerule*[.5pc]{.}\;\thecontentspage} % Formatting of the filler to the right of the heading and the page number

%----------------------------------------------------------------------------------------
%	MINI TABLE OF CONTENTS IN PART HEADS
%----------------------------------------------------------------------------------------

% Chapter text styling
\titlecontents{lchapter}
	[0em] % Left indentation
	{\addvspace{15pt}\large\sffamily\bfseries} % Spacing and font options for chapters
	{\color{azure}\contentslabel[\Large\thecontentslabel]{1.25cm}\color{azure}} % Chapter number
	{}  
	{\color{azure}\normalsize\sffamily\bfseries\;\titlerule*[.5pc]{.}\;\thecontentspage} % Page number

% Section text styling
\titlecontents{lsection}
	[0em] % Left indentation
	{\sffamily\small} % Spacing and font options for sections
	{\contentslabel[\thecontentslabel]{1.25cm}} % Section number
	{}
	{}

% Subsection text styling (note these aren't shown by default, display them by searchings this file for tocdepth and reading the commented text)
\titlecontents{lsubsection}
	[.5em] % Left indentation
	{\sffamily\footnotesize} % Spacing and font options for subsections
	{\contentslabel[\thecontentslabel]{1.25cm}}
	{}
	{}

%----------------------------------------------------------------------------------------
%	HEADERS AND FOOTERS
%----------------------------------------------------------------------------------------

\usepackage{fancyhdr} % Required for header and footer configuration

\pagestyle{fancy} % Enable the custom headers and footers

\renewcommand{\chaptermark}[1]{\markboth{\sffamily\normalsize\bfseries\chaptername\ \thechapter.\ #1}{}} % Styling for the current chapter in the header
\renewcommand{\sectionmark}[1]{\markright{\sffamily\normalsize\thesection\hspace{5pt}#1}{}} % Styling for the current section in the header

\fancyhf{} % Clear default headers and footers
\fancyhead[LE,RO]{\sffamily\normalsize\thepage} % Styling for the page number in the header
\fancyhead[LO]{\rightmark} % Print the nearest section name on the left side of odd pages
\fancyhead[RE]{\leftmark} % Print the current chapter name on the right side of even pages
%\fancyfoot[C]{\thepage} % Uncomment to include a footer

\renewcommand{\headrulewidth}{0.5pt} % Thickness of the rule under the header

\fancypagestyle{plain}{% Style for when a plain pagestyle is specified
	\fancyhead{}\renewcommand{\headrulewidth}{0pt}%
}

% Removes the header from odd empty pages at the end of chapters
\makeatletter
\renewcommand{\cleardoublepage}{
\clearpage\ifodd\c@page\else
\hbox{}
\vspace*{\fill}
\thispagestyle{empty}
\newpage
\fi}

%----------------------------------------------------------------------------------------
%	THEOREM STYLES
%----------------------------------------------------------------------------------------

\usepackage{amsmath,amsfonts,amssymb,amsthm} % For math equations, theorems, symbols, etc

\newcommand{\intoo}[2]{\mathopen{]}#1\,;#2\mathclose{[}}
\newcommand{\ud}{\mathop{\mathrm{{}d}}\mathopen{}}
\newcommand{\intff}[2]{\mathopen{[}#1\,;#2\mathclose{]}}
\renewcommand{\qedsymbol}{$\blacksquare$}
\newtheorem{notation}{Notation}[chapter]

% Boxed/framed environments
\newtheoremstyle{ocrenumbox}% Theorem style name
{0pt}% Space above
{0pt}% Space below
{\normalfont}% Body font
{}% Indent amount
{\small\bf\sffamily\color{azure}}% Theorem head font
{\;}% Punctuation after theorem head
{0.25em}% Space after theorem head
{\small\sffamily\color{azure}\thmname{#1}\nobreakspace\thmnumber{\@ifnotempty{#1}{}\@upn{#2}}% Theorem text (e.g. Theorem 2.1)
\thmnote{\nobreakspace\the\thm@notefont\sffamily\bfseries\color{black}---\nobreakspace#3.}} % Optional theorem note

\newtheoremstyle{blacknumex}% Theorem style name
{5pt}% Space above
{5pt}% Space below
{\normalfont}% Body font
{} % Indent amount
{\small\bf\sffamily}% Theorem head font
{\;}% Punctuation after theorem head
{0.25em}% Space after theorem head
{\small\sffamily{\tiny\ensuremath{\blacksquare}}\nobreakspace\thmname{#1}\nobreakspace\thmnumber{\@ifnotempty{#1}{}\@upn{#2}}% Theorem text (e.g. Theorem 2.1)
\thmnote{\nobreakspace\the\thm@notefont\sffamily\bfseries---\nobreakspace#3.}}% Optional theorem note

\newtheoremstyle{blacknumbox} % Theorem style name
{0pt}% Space above
{0pt}% Space below
{\normalfont}% Body font
{}% Indent amount
{\small\bf\sffamily}% Theorem head font
{\;}% Punctuation after theorem head
{0.25em}% Space after theorem head
{\small\sffamily\thmname{#1}\nobreakspace\thmnumber{\@ifnotempty{#1}{}\@upn{#2}}% Theorem text (e.g. Theorem 2.1)
\thmnote{\nobreakspace\the\thm@notefont\sffamily\bfseries---\nobreakspace#3.}}% Optional theorem note

% Non-boxed/non-framed environments
\newtheoremstyle{ocrenum}% Theorem style name
{5pt}% Space above
{5pt}% Space below
{\normalfont}% Body font
{}% Indent amount
{\small\bf\sffamily\color{azure}}% Theorem head font
{\;}% Punctuation after theorem head
{0.25em}% Space after theorem head
{\small\sffamily\color{azure}\thmname{#1}\nobreakspace\thmnumber{\@ifnotempty{#1}{}\@upn{#2}}% Theorem text (e.g. Theorem 2.1)
\thmnote{\nobreakspace\the\thm@notefont\sffamily\bfseries\color{black}---\nobreakspace#3.}} % Optional theorem note
\makeatother

% Defines the theorem text style for each type of theorem to one of the three styles above
\newcounter{dummy} 
\numberwithin{dummy}{section}
\theoremstyle{ocrenumbox}
\newtheorem{theoremeT}[dummy]{Theorem}
\newtheorem{problem}{Problem}[chapter]
\newtheorem{exerciseT}{Exercise}[chapter]
\theoremstyle{blacknumex}
\newtheorem{exampleT}{Example}[chapter]
\theoremstyle{blacknumbox}
\newtheorem{vocabulary}{Vocabulary}[chapter]
\newtheorem{definitionT}{Definition}[section]
\newtheorem{corollaryT}[dummy]{Corollary}
\theoremstyle{ocrenum}
\newtheorem{proposition}[dummy]{Proposition}

%----------------------------------------------------------------------------------------
%	DEFINITION OF COLORED BOXES
%----------------------------------------------------------------------------------------

\RequirePackage[framemethod=default]{mdframed} % Required for creating the theorem, definition, exercise and corollary boxes

% Theorem box
\newmdenv[skipabove=7pt,
skipbelow=7pt,
backgroundcolor=black!5,
linecolor=azure,
innerleftmargin=5pt,
innerrightmargin=5pt,
innertopmargin=5pt,
leftmargin=0cm,
rightmargin=0cm,
innerbottommargin=5pt]{tBox}

% Exercise box	  
\newmdenv[skipabove=7pt,
skipbelow=7pt,
rightline=false,
leftline=true,
topline=false,
bottomline=false,
backgroundcolor=azure!10,
linecolor=azure,
innerleftmargin=5pt,
innerrightmargin=5pt,
innertopmargin=5pt,
innerbottommargin=5pt,
leftmargin=0cm,
rightmargin=0cm,
linewidth=4pt]{eBox}	

% Definition box
\newmdenv[skipabove=7pt,
skipbelow=7pt,
rightline=false,
leftline=true,
topline=false,
bottomline=false,
linecolor=azure,
innerleftmargin=5pt,
innerrightmargin=5pt,
innertopmargin=0pt,
leftmargin=0cm,
rightmargin=0cm,
linewidth=4pt,
innerbottommargin=0pt]{dBox}	

% Corollary box
\newmdenv[skipabove=7pt,
skipbelow=7pt,
rightline=false,
leftline=true,
topline=false,
bottomline=false,
linecolor=gray,
backgroundcolor=black!5,
innerleftmargin=5pt,
innerrightmargin=5pt,
innertopmargin=5pt,
leftmargin=0cm,
rightmargin=0cm,
linewidth=4pt,
innerbottommargin=5pt]{cBox}

% Creates an environment for each type of theorem and assigns it a theorem text style from the "Theorem Styles" section above and a colored box from above
\newenvironment{theorem}{\begin{tBox}\begin{theoremeT}}{\end{theoremeT}\end{tBox}}
\newenvironment{exercise}{\begin{eBox}\begin{exerciseT}}{\hfill{\color{azure}\tiny\ensuremath{\blacksquare}}\end{exerciseT}\end{eBox}}				  
\newenvironment{definition}{\begin{dBox}\begin{definitionT}}{\end{definitionT}\end{dBox}}	
\newenvironment{example}{\begin{exampleT}}{\hfill{\tiny\ensuremath{\blacksquare}}\end{exampleT}}		
\newenvironment{corollary}{\begin{cBox}\begin{corollaryT}}{\end{corollaryT}\end{cBox}}	

%----------------------------------------------------------------------------------------
%	REMARK ENVIRONMENT
%----------------------------------------------------------------------------------------

\newenvironment{remark}{\par\vspace{10pt}\small % Vertical white space above the remark and smaller font size
\begin{list}{}{
\leftmargin=35pt % Indentation on the left
\rightmargin=25pt}\item\ignorespaces % Indentation on the right
\makebox[-2.5pt]{\begin{tikzpicture}[overlay]
\node[draw=azure!60,line width=1pt,circle,fill=azure!25,font=\sffamily\bfseries,inner sep=2pt,outer sep=0pt] at (-15pt,0pt){\textcolor{azure}{R}};\end{tikzpicture}} % Orange R in a circle
\advance\baselineskip -1pt}{\end{list}\vskip5pt} % Tighter line spacing and white space after remark

%----------------------------------------------------------------------------------------
%	SECTION NUMBERING IN THE MARGIN
%----------------------------------------------------------------------------------------

\makeatletter
\renewcommand{\@seccntformat}[1]{\llap{\textcolor{azure}{\csname the#1\endcsname}\hspace{1em}}}                    
\renewcommand{\section}{\@startsection{section}{1}{\z@}
{-4ex \@plus -1ex \@minus -.4ex}
{1ex \@plus.2ex }
{\normalfont\large\sffamily\bfseries}}
\renewcommand{\subsection}{\@startsection {subsection}{2}{\z@}
{-3ex \@plus -0.1ex \@minus -.4ex}
{0.5ex \@plus.2ex }
{\normalfont\sffamily\bfseries}}
\renewcommand{\subsubsection}{\@startsection {subsubsection}{3}{\z@}
{-2ex \@plus -0.1ex \@minus -.2ex}
{.2ex \@plus.2ex }
{\normalfont\small\sffamily\bfseries}}                        
\renewcommand\paragraph{\@startsection{paragraph}{4}{\z@}
{-2ex \@plus-.2ex \@minus .2ex}
{.1ex}
{\normalfont\small\sffamily\bfseries}}

%----------------------------------------------------------------------------------------
%	PART HEADINGS
%----------------------------------------------------------------------------------------

% Numbered part in the table of contents
\newcommand{\@mypartnumtocformat}[2]{%
	\setlength\fboxsep{0pt}%
	\noindent\colorbox{azure!20}{\strut\parbox[c][.7cm]{\ecart}{\color{azure!70}\Large\sffamily\bfseries\centering#1}}\hskip\esp\colorbox{azure!40}{\strut\parbox[c][.7cm]{\linewidth-\ecart-\esp}{\Large\sffamily\centering#2}}%
}

% Unnumbered part in the table of contents
\newcommand{\@myparttocformat}[1]{%
	\setlength\fboxsep{0pt}%
	\noindent\colorbox{azure!40}{\strut\parbox[c][.7cm]{\linewidth}{\Large\sffamily\centering#1}}%
}

\newlength\esp
\setlength\esp{4pt}
\newlength\ecart
\setlength\ecart{1.2cm-\esp}
\newcommand{\thepartimage}{}%
\newcommand{\partimage}[1]{\renewcommand{\thepartimage}{#1}}%
\def\@part[#1]#2{%
\ifnum \c@secnumdepth >-2\relax%
\refstepcounter{part}%
\addcontentsline{toc}{part}{\texorpdfstring{\protect\@mypartnumtocformat{\thepart}{#1}}{\partname~\thepart\ ---\ #1}}
\else%
\addcontentsline{toc}{part}{\texorpdfstring{\protect\@myparttocformat{#1}}{#1}}%
\fi%
\startcontents%
\markboth{}{}%
{\thispagestyle{empty}%
\begin{tikzpicture}[remember picture,overlay]%
\node at (current page.north west){\begin{tikzpicture}[remember picture,overlay]%	
\fill[azure!20](0cm,0cm) rectangle (\paperwidth,-\paperheight);
\node[anchor=north] at (4cm,-3.25cm){\color{azure!40}\fontsize{220}{100}\sffamily\bfseries\thepart}; 
\node[anchor=south east] at (\paperwidth-1cm,-\paperheight+1cm){\parbox[t][][t]{8.5cm}{
\printcontents{l}{0}{\setcounter{tocdepth}{1}}% The depth to which the Part mini table of contents displays headings; 0 for chapters only, 1 for chapters and sections and 2 for chapters, sections and subsections
}};
\node[anchor=north east] at (\paperwidth-1.5cm,-3.25cm){\parbox[t][][t]{15cm}{\strut\raggedleft\color{white}\fontsize{30}{30}\sffamily\bfseries#2}};
\end{tikzpicture}};
\end{tikzpicture}}%
\@endpart}
\def\@spart#1{%
\startcontents%
\phantomsection
{\thispagestyle{empty}%
\begin{tikzpicture}[remember picture,overlay]%
\node at (current page.north west){\begin{tikzpicture}[remember picture,overlay]%	
\fill[azure!20](0cm,0cm) rectangle (\paperwidth,-\paperheight);
\node[anchor=north east] at (\paperwidth-1.5cm,-3.25cm){\parbox[t][][t]{15cm}{\strut\raggedleft\color{white}\fontsize{30}{30}\sffamily\bfseries#1}};
\end{tikzpicture}};
\end{tikzpicture}}
\addcontentsline{toc}{part}{\texorpdfstring{%
\setlength\fboxsep{0pt}%
\noindent\protect\colorbox{azure!40}{\strut\protect\parbox[c][.7cm]{\linewidth}{\Large\sffamily\protect\centering #1\quad\mbox{}}}}{#1}}%
\@endpart}
\def\@endpart{\vfil\newpage
\if@twoside
\if@openright
\null
\thispagestyle{empty}%
\newpage
\fi
\fi
\if@tempswa
\twocolumn
\fi}

%----------------------------------------------------------------------------------------
%	CHAPTER HEADINGS
%----------------------------------------------------------------------------------------

% A switch to conditionally include a picture, implemented by Christian Hupfer
\newif\ifusechapterimage
\usechapterimagefalse
\newcommand{\thechapterimage}{}%
\newcommand{\chapterimage}[1]{\ifusechapterimage\renewcommand{\thechapterimage}{#1}\fi}%
\newcommand{\autodot}{.}
\def\@makechapterhead#1{%
{\parindent \z@ \raggedright \normalfont
\ifnum \c@secnumdepth >\m@ne
\if@mainmatter
\begin{tikzpicture}[remember picture,overlay]
\node at (current page.north west)
{\begin{tikzpicture}[remember picture,overlay]
\node[anchor=north west,inner sep=0pt] at (0,0) {\ifusechapterimage\includegraphics[width=\paperwidth]{\thechapterimage}\fi};
\draw[anchor=west] (\Gm@lmargin,-9cm) node [line width=2pt,rounded corners=15pt,draw=azure,fill=white,fill opacity=0.5,inner sep=15pt]{\strut\makebox[22cm]{}};
\draw[anchor=west] (\Gm@lmargin+.3cm,-9cm) node {\huge\sffamily\bfseries\color{black}\thechapter\autodot~#1\strut};
\end{tikzpicture}};
\end{tikzpicture}
\else
\begin{tikzpicture}[remember picture,overlay]
\node at (current page.north west)
{\begin{tikzpicture}[remember picture,overlay]
\node[anchor=north west,inner sep=0pt] at (0,0) {\ifusechapterimage\includegraphics[width=\paperwidth]{\thechapterimage}\fi};
\draw[anchor=west] (\Gm@lmargin,-9cm) node [line width=2pt,rounded corners=15pt,draw=azure,fill=white,fill opacity=0.5,inner sep=15pt]{\strut\makebox[22cm]{}};
\draw[anchor=west] (\Gm@lmargin+.3cm,-9cm) node {\huge\sffamily\bfseries\color{black}#1\strut};
\end{tikzpicture}};
\end{tikzpicture}
\fi\fi\par\vspace*{270\p@}}}

%-------------------------------------------

\def\@makeschapterhead#1{%
\begin{tikzpicture}[remember picture,overlay]
\node at (current page.north west)
{\begin{tikzpicture}[remember picture,overlay]
\node[anchor=north west,inner sep=0pt] at (0,0) {\ifusechapterimage\includegraphics[width=\paperwidth]{\thechapterimage}\fi};
\draw[anchor=west] (\Gm@lmargin,-9cm) node [line width=2pt,rounded corners=15pt,draw=azure,fill=white,fill opacity=0.5,inner sep=15pt]{\strut\makebox[22cm]{}};
\draw[anchor=west] (\Gm@lmargin+.3cm,-9cm) node {\huge\sffamily\bfseries\color{black}#1\strut};
\end{tikzpicture}};
\end{tikzpicture}
\par\vspace*{270\p@}}
\makeatother

%----------------------------------------------------------------------------------------
%	LINKS
%----------------------------------------------------------------------------------------

\usepackage{hyperref}
\hypersetup{hidelinks,backref=true,pagebackref=true,hyperindex=true,colorlinks=false,breaklinks=true,urlcolor=azure,bookmarks=true,bookmarksopen=false}

\usepackage{bookmark}
\bookmarksetup{
open,
numbered,
addtohook={%
\ifnum\bookmarkget{level}=0 % chapter
\bookmarksetup{bold}%
\fi
\ifnum\bookmarkget{level}=-1 % part
\bookmarksetup{color=azure,bold}%
\fi
}
}
 % Insert the commands.tex file which contains the majority of the structure behind the template

%\hypersetup{pdftitle={Title},pdfauthor={Author}} % Uncomment and fill out to include PDF metadata for the author and title of the book

%----------------------------------------------------------------------------------------

\begin{document}

%----------------------------------------------------------------------------------------
%	TITLE PAGE
%----------------------------------------------------------------------------------------

\begingroup
\thispagestyle{empty} % Suppress headers and footers on the title page
\begin{tikzpicture}[remember picture,overlay]
%\node[inner sep=0pt] (background) at (current page.center) {\includegraphics[width=\paperwidth]{background.pdf}};
\draw (current page.center) node [fill=azure!30!white,fill opacity=0.6,text opacity=1,inner sep=1cm]{\Huge\centering\bfseries\sffamily\parbox[c][][t]{\paperwidth}{\centering M4997: Parallel Computaitonal Mathematics\\[15pt] % Book title
{\Large Fall 2020}\\[20pt] % Subtitle
{\huge Dr. Patrick Diehl}}}; % Author name
\end{tikzpicture}
\vfill
\endgroup

%----------------------------------------------------------------------------------------
%	COPYRIGHT PAGE
%----------------------------------------------------------------------------------------

\newpage
~\vfill
\thispagestyle{empty}

\noindent Copyright \copyright\ 2020 Patrick Diehl\orcid{0000-0003-3922-8419}\\ % Copyright notice

%\noindent \textsc{Published by Publisher}\\ % Publisher

\noindent \url{https://www.cct.lsu.edu/~pdiehl/teaching/2020/4997/}\\ % URL

\noindent \doclicenseThis  

\noindent \textit{Edition, Fall 2020} % Printing/edition date

%----------------------------------------------------------------------------------------
%	TABLE OF CONTENTS
%----------------------------------------------------------------------------------------

%\usechapterimagefalse % If you don't want to include a chapter image, use this to toggle images off - it can be enabled later with \usechapterimagetrue

\chapterimage{chapter_head_1.pdf} % Table of contents heading image

\pagestyle{empty} % Disable headers and footers for the following pages

\tableofcontents % Print the table of contents itself

\cleardoublepage % Forces the first chapter to start on an odd page so it's on the right side of the book

\pagestyle{fancy} % Enable headers and footers again

%----------------------------------------------------------------------------------------
%	PART
%----------------------------------------------------------------------------------------

\part{Introduction: C++ and the C++ standard template library}

%----------------------------------------------------------------------------------------
%	CHAPTER 1
%----------------------------------------------------------------------------------------

\chapterimage{chapter_head_2.pdf} % Chapter heading image

%----------------------------------------------------------------------------------------
\chapter{Introduction C++}
%----------------------------------------------------------------------------------------

%----------------------------------------------------------------------------------------
\section{History of C and C++}
%----------------------------------------------------------------------------------------

%----------------------------------------------------------------------------------------
\section{Getting started with C++}
%----------------------------------------------------------------------------------------
To begin with C++ programming, we look at a simple C++ program, the so-called ``Hello World'' example which most programming language start with. Listing~\ref{code:hello:world} shows this program and the first line in green a comment is shown. A  single-line comments always starts with \lstinline[language=C++]|//| and is used to explain the functionality of the program or the next lines of codes. It is also possible to use multi-line comments\endnote{\url{https://en.cppreference.com/w/cpp/comment}} by using \lstinline[language=C++]{/* */}. Comments are important to understand the program, especially if the code is shared with other collaborators. Fore more details we refer to~\cite{kernighan1974elements}.

The second line starts with a so-called include directive\endnote{\url{https://en.cppreference.com/w/cpp/preprocessor/include}} \lstinline[language=C++]{#include <iostream>}. This include directive is needed to include functionality of the C++ standard library, see Chpater~\ref{chapter:stl}, in our case we include the \lstinline|iostream| header to print the Hello World to the terminal, see Line 6.

The fourth line \lstinline[language=C++]{int main()} starts with the so-called Main function\endnote{\url{https://en.cppreference.com/w/cpp/language/main_function}} which is the entry point of the program and all code lines within are executed sequentially one by one. Every C++ which will be compiled to an executable needs exact one function called \lstinline[language=C++]|main| which have a integer \lstinline[language=C++]{int} as its return type. On most operation systems a return value of zero means that the program executed successfully and any other value indicates an failure. The second last line \lstinline[]|return| is the so--called return statement\endnote{\url{https://en.cppreference.com/w/cpp/language/return}} which has to match the return type in front of the \lstinline[language=C++]{int main()}.


\lstinputlisting[language=C++,caption={A simple C++ program, the so-called ``Hello World'' example, which most languages start with.\label{code:hello:world}},float,floatplacement=h]{ParallelComputationMathExamples/chapter2/lecture1-main.cpp}

Once we have written the program, we have to compile the C++ code into an executable to run the code and print ``Hello world'' to the terminal. Note that there are plenty of C++ compilers\endnote{\url{https://en.wikipedia.org/wiki/List_of_compilers\#C++_compilers}} available, however this book will use the GNU Compiler Collection (GCC) for all examples. Line 1 in Listing~\ref{code:hello:world:compile} shows how to compile the file \lstinline[language=bash]|lecture1-1.cpp|, which contains the C++ code in Listing~\ref{code:hello:world}, to an executable. The GCC provides the \lstinline[language=bash]{g++} compiler to compile C++ code and the \lstinline[language=bash]{gcc} compiler to compile C code. As the first argument to the \lstinline[language=bash]{g++} compiler the file name of the C++ is provided and with the \lstinline[language=bash]|-o| option the name of the executable is specified. To run the generated executable, we type \lstinline[language=bash]|./lecture-1-1| in the terminal. Note for the basic usage of the Linux terminal we refer to~\cite{newham2005learning,robbins2016bash}.

\lstinputlisting[language=bash,caption={Compilation and execution of the C++ program.\label{code:hello:world:compile}},float,floatplacement=h,firstline=2, lastline=3]{ParallelComputationMathExamples/chapter2/lecture1-main.sh}

\begin{exercise}
Download the example program\endnote{\url{https://github.com/diehlpkteaching/ParallelComputationMathExamples}} from GitHub and compile it with your favorite C++ compiler. After you ran the example you could try to modify it, for example you could print a different text or add a second line to the output.
\end{exercise}


%----------------------------------------------------------------------------------------
\subsection{Fundamental data types}
%----------------------------------------------------------------------------------------
\index{data types!fundamental}
In this section the fundamental data types\endnote{\url{https://en.cppreference.com/w/cpp/language/types}} provided by the C++ language are introduced. First, the numeric data types are introduced. To represent natural numbers $\mathbb{N}=\{0,1,2,\ldots \}$ the \lstinline[language=C++]|unsigned int| data type is available. To represent integer numbers $\mathbb{Z}=\{\ldots,-2,-1,0,1,2,\ldots \}$ the \lstinline[language=C++]|int| data type is available. For all these data types following options: \lstinline[language=C++]|short|, \lstinline[language=C++]|long|, and \lstinline[language=C++]|long long| are available. For more details about the binary numeral system we refer to~\cite{gilli1965binary}. To represent real numbers $\mathbb{R}$ the \lstinline[language=C++]|float| data type and \lstinline[language=C++]|double| data type are available. Fore more details about the IEEE 474 standard how floating point numbers are represented in the computer we refer to~\cite{4610935,goldberg1991every}. Table~\ref{chapter2:table:datatypes} summarizes all the available numeric data types and there ranges. The next section shows how to get the range of the IEEE 474 standard how floating point numbers.

\begin{table}[h]
\centering
\begin{tabular}{lccc}
\toprule
Data type & Size (Bytes) & Min & Max \\\midrule
\multicolumn{4}{c}{Natural numbers $\mathbb{N}$ }\\\midrule
\lstinline[language=C++]|unsigned short int| & 2 & 0 & 65,535  \\ 
\lstinline[language=C++]|unsigned int| & 4 & 0 & 4,294,967,295 \\ 
\lstinline[language=C++]|unsigned long int| & 4 & 0 & 4,294,967,295 \\ 
\lstinline[language=C++]|unsigned long long int| & 8 & 0 & 8,446,744,073,709,551,615 \\ \midrule
\multicolumn{4}{c}{Integer numbers $\mathbb{Z}$ }\\\midrule
\lstinline[language=C++]|short int| & 2 & -32,768 & 32,768 \\
\lstinline[language=C++]|int| & 4 & -2,147,483,648 & 2,147,483,648 \\
\lstinline[language=C++]|long long int| & 8 & $-2^{63}$ & $2^{63}-1$ \\\midrule
\multicolumn{4}{c}{Real numbers $\mathbb{R}$ }\\\midrule
\lstinline[language=C++]|float| & 4 &  &  \\
\lstinline[language=C++]|double| & 8 &  &  \\
\bottomrule
\end{tabular} 
\caption{Overview of the fundamental numeric data types.}
\label{chapter2:table:datatypes}
\end{table}

To represent a boolean value $\mathbf{B}=\{0,1\}$ the \lstinline[language=C++]|bool| data type which has either one of the two values \lstinline[language=C++]|true| or \lstinline[language=C++]|false|. Note that the C++ STL offers \lstinline|std::complex|\endnote{\url{https://en.cppreference.com/w/cpp/numeric/complex}} for complex number $\mathbb{C}$, however, this one is not within the fundamental data types.

\subsection{Numerical limits}


%----------------------------------------------------------------------------------------
\subsection{Statements and flow control}
%----------------------------------------------------------------------------------------

%----------------------------------------------------------------------------------------
\subsubsection{Iteration statements}
%----------------------------------------------------------------------------------------
\index{statements!iteration} For some applications, we have to repeat some operations multiple times. The C++ language provides the so-called iteration statements. There are two iteration statements, the \lstinline[language=C++]|for| loop and the \lstinline[language=C++]|while| loop, respectively. Let us look how to compute the sum of the numbers from one up ton $n$
\begin{align}
r = \sum\limits_{i=1}^n i\text{.}
\end{align}
The first option is using a \lstinline[language=C++]|for| loop statement\endnote{\url{https://en.cppreference.com/w/cpp/language/for}}\index{loop!for statement} which shown in Listing~\ref{code:for:loop}. Line 9 shows the \lstinline[language=C++]|for| loop statement with its three arguments. First, the so--called loop variable \lstinline[language=C++]{size_t i = 0} which is initialized to zero. Note that the loop variable is only defined within the loop's body (The part between the curly braces). Second, the so--called condition statement \lstinline[language=C++]{i < n} which means that the loop body is repeated until the variable $i$ is larger than $n$. The third statement manipulates the loop variable, in our case the loop variable is incremented by one after the loop body was executed once. Note that we use the \lstinline[language=C++]|for| loop statement, if we know how often we have to loop through the loop body in advance.


\lstinputlisting[language=C++,caption={Computation of the sum from one up to $n$ using the for loop statement. \label{code:for:loop}},float,floatplacement=h]{ParallelComputationMathExamples/chapter2/lecture1-for.cpp}

The second option is using a \lstinline[language=C++]|while| loop statement\endnote{\url{https://en.cppreference.com/w/cpp/language/while}}\index{loop!while statement} which shown in Listing~\ref{code:while:loop}. Line 10 shows the \lstinline[language=C++]|while| loop statement with its one argument. This is the so-called condition statement  \lstinline[language=C++]{i < n} which means that the loop body is repeated until the variable $i$ is larger than $n$. Note in the previous example we had three arguments. Here, the loop variable is declared before the loop in Line 9 and the third statement is  in Line 13 where the loop variable is incremented by one in each iteration. Note that we use the \lstinline[language=C++]|while| loop statement, if we do not know the amount of iterations in advance. This example has shown that we can write every \lstinline[language=C++]|for| loop statement as a \lstinline[language=C++]|while| loop statement.   


\lstinputlisting[language=C++,caption={Computation of the sum from one up to $n$ using the while loop statement..\label{code:while:loop}},float,floatplacement=h]{ParallelComputationMathExamples/chapter2/lecture1-while.cpp}

\begin{exercise}
Define in your own words in which case you should use a \lstinline[language=C++]|for| loop statement and a \lstinline[language=C++]|while| loop statement. 
\end{exercise}

%----------------------------------------------------------------------------------------
\subsubsection{Selection statements}
%----------------------------------------------------------------------------------------
\index{statements!selection}

\begin{align}
r = \sum\limits_{i=1}^n f(i) \text{  with  } f(i) = 
\begin{cases}
i, \text{ if } i \text{ is even} \\
i^2, \text{ else}
\end{cases}
\end{align}


%----------------------------------------------------------------------------------------
\subsubsection{Operators}
%----------------------------------------------------------------------------------------
\index{Operators}
For the example in Listing~\ref{code:for:loop} we have seen the operator \lstinline[language=C++]|i<n| which is a so--called comparison operator. Next to the comparison operators, C++ language has following operators\endnote{\url{https://en.cppreference.com/w/cpp/language/operator_precedence}}:
\begin{itemize}
\item Comparison operators, see Table~\ref{sec:1:tab:operator:comp},
\item Arithmetic operators, see Table~\ref{sec:1:tab:operator:arithmetic},
\item Logical operators, see Table~\ref{sec:1:tab:operator:logical}, and
\item Assignment operators, see Table~\ref{sec:1:tab:operator:assign},
\end{itemize}
 logical operators, arithmetic, and assignment.

\begin{table}[h]
\centering
\begin{tabular}{clll}
\toprule
Operator & Name  & Example \\ 
\midrule
\lstinline|==| & Equal to & \lstinline|x==y|\\ 
\lstinline|!=| & Not equal & \lstinline|x!=y|\\ 
\lstinline|>| & Greater than & \lstinline|x > y|\\ 
\lstinline|<| & Less than & \lstinline|x < y|\\ 
\lstinline|>=| & Greater than or equal & \lstinline|x >= y|\\ 
\lstinline|<=| & Less than or equal & \lstinline|x <= y|\\ 
\bottomrule 
\end{tabular} 
\caption{Comparison operators}
\label{sec:1:tab:operator:comp}
\index{Operators!comparison}
\end{table}

\begin{table}[h]
\centering
\begin{tabular}{clll}
\toprule
Operator & Name & Description & Example \\ 
\midrule
\lstinline|+| & Addition & Computes the sum of two values & $2+2=4$ \\ 
\lstinline|-| & Subtraction  & Computes the difference of two values & $5-3=2$ \\ 
\lstinline|/| & Division & Divides two values & $6/2=3$ \\ 
\lstinline|*| & Multiplication & Multiplies two values & $2\times2=4$ \\ 
\lstinline|%| & Modulo &  	Returns the division remainder & \lstinline|2%1=0| \\ 
\lstinline|++| & Increments & Add plus one to the value & \lstinline|1++=2|\\ 
\lstinline|--| & Decrements & Subtract one of the value & \lstinline|1--=0|\\ 
\bottomrule 
\end{tabular} 
\caption{Arithmetic operators}
\label{sec:1:tab:operator:arithmetic}
\index{Operators!arithmetic}
\end{table}


\begin{table}[h]
\centering
\begin{tabular}{clll}
\toprule
Operator & Name & Description & Example \\ 
\midrule
\lstinline|&&| & Logical and & Returns \lstinline|true| if both statements are true  & \lstinline| x > 5 && x < 10| \\ 
\lstinline|||| & Logical or  & Returns \lstinline|true| if one statement is true & \lstinline| x > 5 || y < 10| \\ 
\lstinline|!| & Logical not &  Inverse the statement & \lstinline| !(x > 5 && x < 10)| \\ 
\bottomrule 
\end{tabular} 
\caption{Logical operators}
\label{sec:1:tab:operator:logical}
\index{Operators!logical}
\end{table}


\begin{table}[h]
\centering
\begin{tabular}{clll}
\toprule
Operator & Name & Example & Equivalent  \\ 
\midrule
\lstinline|=| & Assignment &   \lstinline| x = 5| &   \lstinline| x = 5 | \\ 
\lstinline|+=| & Plus equal  & \lstinline| x+= 5|  & \lstinline| x = x + 5 | \\ 
\lstinline|-=| & Minus equal & \lstinline| x-= 5|  & \lstinline| x = x - 5 | \\ 
\lstinline|*=| & Multiplication equal &  \lstinline| x*= 5| & \lstinline| x= x * 5| \\ 
\lstinline|/=| & Division equal &  \lstinline| x/= 5| & \lstinline| x= x / 5| \\ 
\lstinline|%=| & Modulo equal &  \lstinline| x%= 5| & \lstinline| x = x % 5| \\ 
\bottomrule 
\end{tabular} 
\caption{Assignment operators}
\label{sec:1:tab:operator:assign}
\index{Operators!assignment}
\end{table}

\begin{exercise}
Write a small C++ program using selection statements and operators to determine if a given year is a lap year. Following logical statements should be implemented: 
\begin{itemize}
	\item   If year is divided by 4 but not by 100, then it is a leap year.
    \item If year is divided by both 100 and 400, then it is a leap year.
    \item If year is divided by 400, then it is a leap year.
    \item And in all other cases, it is not a leap year.
\end{itemize}
\end{exercise}

%----------------------------------------------------------------------------------------
\subsection{Functions}
%----------------------------------------------------------------------------------------

%----------------------------------------------------------------------------------------
\subsection{Structuring source code}
%----------------------------------------------------------------------------------------

%----------------------------------------------------------------------------------------
\subsubsection{Structs}
%----------------------------------------------------------------------------------------

%----------------------------------------------------------------------------------------
\subsection{Classes}
%----------------------------------------------------------------------------------------

\newpage

\theendnotes

\chapter{The C++ Standard Template Library (STL)}
\label{chapter:stl}

\section{Overview of the STL}

\section{Working with strings}

\section{Containers}


\part{Linear algebra and solvers}

\chapter{Linear algebra}

\section{Blaze library}

\chapter{Solvers}


\part{Numerical examples}

\chapter{Numerical examples}

\section{Monte-Carlo methods}

\section{$N$-body problems}

\section{One-dimensional heat equation}

\section{Peridynamics}

\part{Parallel and distributes computing}

\chapter{Parallel computing}

\section{Shared memory}

\section{Distributed memory}

\part{Introduction: HPX}

\chapter{HPX}

\section{Introduction to HPX}

\section{Parallel algorithms}

\section{Asynchronous programming}





%----------------------------------------------------------------------------------------
%	BIBLIOGRAPHY
%----------------------------------------------------------------------------------------

\chapter*{Bibliography}
\addcontentsline{toc}{chapter}{\textcolor{azure}{Bibliography}} % Add a Bibliography heading to the table of contents

%------------------------------------------------

\section*{Articles}
\addcontentsline{toc}{section}{Articles}
\printbibliography[heading=bibempty,type=article]

%------------------------------------------------

\section*{Books}
\addcontentsline{toc}{section}{Books}
\printbibliography[heading=bibempty,type=book]

%----------------------------------------------------------------------------------------
%	INDEX
%----------------------------------------------------------------------------------------

\cleardoublepage % Make sure the index starts on an odd (right side) page
\phantomsection
\setlength{\columnsep}{0.75cm} % Space between the 2 columns of the index
\addcontentsline{toc}{chapter}{\textcolor{azure}{Index}} % Add an Index heading to the table of contents
\printindex % Output the index

%----------------------------------------------------------------------------------------

%----------------------------------------------------------------------------------------
%	List of Listings
%----------------------------------------------------------------------------------------

\cleardoublepage % Make sure the index starts on an odd (right side) page
\phantomsection
\addcontentsline{toc}{chapter}{\textcolor{azure}{List of Listings}}
\lstlistoflistings


\end{document}
